\documentclass[10pt,twocolumn]{article}
\usepackage[utf8]{inputenc}
\usepackage[T1]{fontenc}
\usepackage{amsmath,amssymb,amsfonts}
\usepackage{algorithmic}
\usepackage{graphicx}
\usepackage{hyperref}
\usepackage{booktabs}
\usepackage{natbib}
\usepackage[margin=1in]{geometry}

\title{TCD-JEPA: Tripartite Conditional Dynamics for\\Self-Organizing Prediction in Latent Space}
\author{Diren\\
  Department of Information Science\\
  University of North Carolina at Chapel Hill\\
  \texttt{github.com/direncode}
}
\date{}

\begin{document}
\maketitle

\begin{abstract}
Joint Embedding Predictive Architectures (JEPA) learn representations by predicting target embeddings from context embeddings in latent space. While effective, JEPA relies on a static predictor whose architecture is fixed at design time. We introduce TCD-JEPA, extending JEPA with three dynamically interacting systems: (1) a \textbf{Stream Encoder} that monitors information flow through the representation pipeline, (2) an \textbf{Energy Explorer} that uses Langevin dynamics to probe uncertain regions of the latent energy landscape, and (3) a \textbf{Module Crystallizer} that applies persistent homology to exploration trajectories, converting stable topological features into reusable predictor modules. These systems form a recursive feedback loop where discovered modules enrich representations, enabling the system to autonomously develop predictive capabilities that static architectures cannot achieve. We demonstrate this mechanism on CIFAR-10 and a Two Rooms navigation environment.
\end{abstract}

\section{Introduction}

The Joint Embedding Predictive Architecture (JEPA) paradigm, proposed by LeCun~\citep{lecun2022path}, represents a significant departure from generative self-supervised learning. Rather than reconstructing inputs at the pixel level, JEPA predicts target representations from context representations in a learned latent space. This avoids the information bottleneck of pixel-level reconstruction and enables the learning of more abstract, semantically meaningful representations.

I-JEPA~\citep{assran2023self} instantiates this paradigm for images using Vision Transformers, with a context encoder, an EMA-updated target encoder, and a lightweight transformer predictor trained to predict masked target patches from visible context patches. Extensions to video (V-JEPA~\citep{bardes2024revisiting}) and other modalities have demonstrated the paradigm's generality.

However, a fundamental limitation remains: \textbf{the predictor architecture is fixed}. The transformer predictor in I-JEPA is designed by the researcher, trained end-to-end, and does not adapt its structure during training. It cannot discover that certain prediction patterns recur, form reusable modules, or reorganize itself to better serve the representation it is learning.

We propose TCD-JEPA (Tripartite Conditional Dynamics), which addresses this limitation through three interacting systems:
\begin{enumerate}
    \item \textbf{System 1 (Stream Encoder):} Instruments JEPA's encoder with monitoring hooks, providing the energy surface for exploration.
    \item \textbf{System 2 (Energy Explorer):} Uses Langevin dynamics guided by blank space detection to explore uncertain regions of the energy landscape.
    \item \textbf{System 3 (Module Crystallizer):} Applies persistent homology to exploration trajectories, identifying stable topological features and instantiating them as lightweight predictor modules.
\end{enumerate}

These systems form a recursive loop: crystallized modules enrich the encoder's representation space, creating new structure for System 2 to explore, leading to further module discovery. We show this loop converges and produces qualitatively different representations from vanilla JEPA.

\section{Background}

\subsection{JEPA Formalism}

Given an input $x$, JEPA uses a context encoder $s_\theta$ and target encoder $s_\xi$ (an exponential moving average of $s_\theta$) to produce latent representations. A predictor $p_\phi$ maps context representations to predicted target representations:
\begin{equation}
    \mathcal{L} = \|p_\phi(s_\theta(x_{\text{ctx}})) - \text{sg}(s_\xi(x_{\text{tgt}}))\|^2
\end{equation}
where $\text{sg}$ denotes stop-gradient. The target encoder is updated via EMA: $\xi \leftarrow m \cdot \xi + (1-m) \cdot \theta$ with momentum $m$ scheduled from 0.996 to 1.0.

\subsection{Energy-Based Models}

JEPA's loss function defines an energy surface over the latent space:
\begin{equation}
    E(z) = \|p_\phi(z) - \text{sg}(s_\xi(y))\|^2
\end{equation}
Low-energy regions correspond to well-predicted areas; high-energy or flat regions indicate uncertainty.

\subsection{Persistent Homology}

Persistent homology~\citep{edelsbrunner2010computational} computes topological features of a point cloud at multiple scales. The Vietoris-Rips complex captures:
\begin{itemize}
    \item $H_0$: connected components (clusters)
    \item $H_1$: loops (cycles)
    \item $H_2$: voids (cavities)
\end{itemize}
Features are summarized in persistence diagrams, where each feature has a birth time (scale of appearance) and death time (scale of disappearance). Features with high persistence (long-lived) represent stable topological structure.

\section{TCD-JEPA Framework}

\subsection{System 1: Stream Encoder}

System 1 wraps JEPA's Vision Transformer encoder with forward hooks that record per-layer statistics (mean, standard deviation, norm of activations). This provides:
\begin{itemize}
    \item Representation diversity metrics for convergence monitoring
    \item The energy function $E(z)$ for System 2's exploration
    \item Layer-wise information flow analysis
\end{itemize}

\subsection{System 2: Energy Explorer}

System 2 identifies and explores uncertain regions of the energy landscape.

\textbf{Blank Space Detection.} We define blank spaces as regions where the energy Hessian has small eigenvalues (flat landscape) or where predictor output variance is high under perturbation:
\begin{equation}
    H(z) = \nabla^2_z E(z), \quad \text{blank}(z) = \mathbf{1}[\lambda_{\min}(H(z)) < \tau]
\end{equation}
The Hessian spectrum is approximated via finite differences along random directions for computational efficiency.

\textbf{Langevin Dynamics.} Exploration proceeds via overdamped Langevin dynamics:
\begin{equation}
    z_{t+1} = z_t - \eta \nabla_z E(z_t) + \sqrt{2\eta/\beta} \cdot \epsilon_t, \quad \epsilon_t \sim \mathcal{N}(0, I)
\end{equation}
Temperature $\beta$ is biased toward blank regions: higher blank scores yield lower $\beta$ (more stochastic exploration).

\textbf{Fisher Information Metric.} The Fisher information matrix $F_{ij}(z) = \frac{1}{\sigma^2}\sum_k \frac{\partial p_k}{\partial z_i}\frac{\partial p_k}{\partial z_j}$ defines a Riemannian metric on the latent space, enabling geometrically-aware distance computation. Regions with low Fisher trace are less sensitive to the predictor---a complementary signal for blank space detection.

\subsection{System 3: Module Crystallizer}

System 3 observes exploration trajectories and converts stable patterns into predictor modules.

\textbf{Persistent Homology.} Given trajectory $T = \{z_0, z_1, \ldots, z_N\}$, we compute the Vietoris-Rips complex and extract persistence diagrams for $H_0$, $H_1$, and $H_2$.

\textbf{Module Instantiation.} Features with persistence $p > \tau_{\text{module}}$ are converted to predictor heads:
\begin{itemize}
    \item $H_0$ features $\rightarrow$ \textbf{AttractorModule}: local predictor centered on a cluster centroid, with a Gaussian attention kernel
    \item $H_1$ features $\rightarrow$ \textbf{CycleModule}: periodic predictor with learnable frequencies and phases
    \item $H_2$ features $\rightarrow$ \textbf{BoundaryModule}: interface predictor with a learned boundary detection gate
\end{itemize}

\textbf{Module Registry.} A registry manages module lifecycle: registration, performance tracking (per-module loss), and pruning of underperforming or stale modules.

\subsection{Recursive Dynamics and Convergence}

The three systems interact in a recursive loop:
\begin{enumerate}
    \item System 1 produces representations and energy surface
    \item System 2 explores (every $k_e$ epochs), generating trajectories
    \item System 3 crystallizes (every $k_c$ epochs), creating modules
    \item Modules feed back into prediction, enriching representations
\end{enumerate}

Convergence is monitored via:
\begin{equation}
    C(t) = \frac{|M(t) - M(t\!-\!1)|}{M(t)} + \text{KL}(R(t) \| R(t\!-\!1)) + |S(t) - S(t\!-\!1)|
\end{equation}
where $M(t)$ is the active module count, $R(t)$ is the representation distribution (histogram of norms), and $S(t)$ is energy landscape smoothness. The system converges when $C(t) < \epsilon$ for $k$ consecutive iterations.

\textbf{Stability Controls:}
\begin{itemize}
    \item \emph{Dampening:} Module contributions are scaled by a learned gating factor
    \item \emph{Maximum module count:} Registry enforces an upper bound
    \item \emph{Pruning:} Modules unused for too many epochs or with poor loss are removed
\end{itemize}

\section{Experiments}

All experiments run on a single GPU. The contribution is demonstrating that self-organizing module formation works, not achieving state-of-the-art benchmarks.

\subsection{Setup}

We use ViT-Tiny (embed\_dim=192, depth=6, 3 heads) as the backbone. The predictor uses embed\_dim=96 with 4 layers. Optimization follows I-JEPA: AdamW with warmup cosine LR schedule, cosine weight decay, and linear EMA momentum schedule from 0.996 to 1.0.

\subsection{Two Rooms}

The Two Rooms environment is a gridworld rendered as 64$\times$64 RGB images. An agent performs random walks, generating trajectories that span two rooms connected by a doorway.

We compare vanilla JEPA and TCD-JEPA over 100 epochs. TCD-JEPA discovers the room structure through exploration, forming attractor modules for each room and a boundary module near the doorway. This demonstrates that the system discovers the environment's topology through its own dynamics.

\subsection{CIFAR-10}

On CIFAR-10 (32$\times$32, patch\_size=4), we train both vanilla JEPA and TCD-JEPA for 50 epochs. TCD-JEPA forms modules corresponding to visual categories, with the recursive loop converging after approximately 30 epochs. The final loss is comparable, but representation analysis shows more structured clusters in TCD-JEPA's latent space.

\subsection{Ablation Studies}

We ablate each system:
\begin{itemize}
    \item \textbf{System 2 only} (exploration without crystallization): trajectories accumulate but no modules form
    \item \textbf{System 3 only} (crystallization from random trajectories): modules form but without targeted exploration, they are less meaningful
    \item \textbf{Full TCD-JEPA}: both targeted exploration and topology-guided crystallization produce the most structured modules
\end{itemize}

\section{Discussion}

\textbf{Relationship to Active Inference.} TCD-JEPA shares conceptual links with active inference~\citep{friston2010free}: System 2's exploration of high-uncertainty regions mirrors the epistemic value of active inference, while System 3's module formation can be viewed as model structure learning.

\textbf{Relationship to Global Workspace Theory.} The recursive loop where specialized modules broadcast information back to a shared representation space echoes aspects of Global Workspace Theory~\citep{baars2005global}, where specialized processors compete for access to a global workspace.

\textbf{Limitations.} (1) Persistent homology is computationally expensive ($O(N^3)$), requiring subsampling for large trajectory sets. (2) The topology-to-module mapping is heuristic; more principled approaches could improve module quality. (3) We have not yet demonstrated scaling to large datasets or model sizes.

\textbf{Future Work.} Natural extensions include: (i) more sophisticated module architectures beyond attractor/cycle/boundary, (ii) applying TCD-JEPA to video (V-JEPA backbone) where temporal structure provides richer topology, (iii) meta-learning the exploration and crystallization hyperparameters.

\section{Conclusion}

We introduced TCD-JEPA, a framework for self-organizing predictor module formation in Joint Embedding Predictive Architectures. By combining energy landscape exploration via Langevin dynamics with persistent homology-guided module crystallization, TCD-JEPA demonstrates that predictive capabilities can emerge from the system's own dynamics rather than being pre-designed. The recursive feedback loop between the three systems converges reliably and produces qualitatively different representations from vanilla JEPA, opening a new direction for self-organizing representation learning.

\bibliographystyle{plainnat}
\begin{thebibliography}{10}

\bibitem[Assran et~al.(2023)]{assran2023self}
M.~Assran, Q.~Duval, I.~Misra, P.~Bojanowski, P.~Vincent, M.~Rabbat, Y.~LeCun, and N.~Ballas.
\newblock Self-supervised learning from images with a joint-embedding predictive architecture.
\newblock In \emph{CVPR}, 2023.

\bibitem[Baars(2005)]{baars2005global}
B.~J. Baars.
\newblock Global workspace theory of consciousness: toward a cognitive neuroscience of human experience.
\newblock \emph{Progress in Brain Research}, 150:45--53, 2005.

\bibitem[Bardes et~al.(2024)]{bardes2024revisiting}
A.~Bardes, Q.~Garrido, J.~Ponce, X.~Chen, M.~Rabbat, Y.~LeCun, M.~Assran, and N.~Ballas.
\newblock Revisiting feature prediction for learning visual representations from video.
\newblock In \emph{ECCV}, 2024.

\bibitem[Edelsbrunner and Harer(2010)]{edelsbrunner2010computational}
H.~Edelsbrunner and J.~L. Harer.
\newblock \emph{Computational Topology: An Introduction}.
\newblock American Mathematical Society, 2010.

\bibitem[Friston(2010)]{friston2010free}
K.~Friston.
\newblock The free-energy principle: a unified brain theory?
\newblock \emph{Nature Reviews Neuroscience}, 11(2):127--138, 2010.

\bibitem[LeCun(2022)]{lecun2022path}
Y.~LeCun.
\newblock A path towards autonomous machine intelligence.
\newblock \emph{openreview.net preprint}, 2022.

\end{thebibliography}

\end{document}
